\documentclass[a4paper, 11pt]{article}
\usepackage[margin=0.75in]{geometry}
\usepackage{graphicx, tabularx, hyperref, amsmath}
\usepackage[siunitx]{circuitikz}
\usepackage{booktabs}
\usepackage{pgfplotstable, pgfplots}
\renewcommand{\arraystretch}{1.2}
\usepackage{xepersian}
\settextfont[Scale=0.9]{Dibaj}
\setdigitfont[Scale=1]{Dibaj}

\begin{document}

\begin{center}
\textbf{{\huge گزارش‌کار آزمایش شمارهٔ
1: 
بررسی قانون اهم
}}

\end{center}

\hrule
\noindent\begin{tabular}{rrl}
اعضای گروه: & ابوالفضل احمدی & 403172283\\
& نام و نام‌خانوادگی & 403001122
\end{tabular}\hfill
\begin{tabular}{rr}
تاریخ انجام آزمایش: & 10/08/1404 \\
دستیار آموزشی: &  \\
\end{tabular}\hfill
\begin{tabular}{rr}
گروه:  &  \\
زیرگروه: \lr{A} &  \\
\end{tabular}
\hrule

\section*{اهداف آزمایش و تئوری آزمایش}

ما در این آزمایش به ﺑﺮرﺳﯽ ﺗﺠﺮﺑﯽ ﻗﺎﻧﻮن اﻫﻢ و ﻣﻄﺎﻟﻌﻪ ﭘﺎراﻣﺘﺮﻫﺎي ﻣﺆﺛﺮ در ﻣﻘﺎوﻣﺖ اﻟﮑﺘﺮﯾﮑﯽ ﯾﮏ ﺳﯿﻢ ﻓﻠﺰي پرداختیم.

\section{بستگی اختلاف پتانسیل دو سر سیم به اندازه جریان الکتریکی عبوری از آن}

\begin{figure}[h!]
\centering
\caption{مدار ساده اندازه‌گیری}
\begin{circuitikz}
  %--- battery on the left ---
  \draw (0,0) to[battery1] (0,3);
  % polarity labels
  %\node at (-0.7,2.3) {$-$};
  %\node at (-0.7,0.7) {$+$};

  % top and bottom rails to the right
  \draw (0,3) -- (2,3);
  \draw (0,0) -- (1,0) to[ammeter] (2,0);

  % vertical voltmeter branch
  \draw (2,3) to[voltmeter] (2,0);

  %--- diverging leads a and b with tap points (red dots) ---
  % style for tap points
  \tikzset{tap/.style={draw=red!60!black, fill=red, circle, inner sep=1.2pt}}

  % upper lead "a"
  \draw
   	(2,3) -- (2.5,3) -- (3,2) -- (8,1.5);
  \draw
   	(2,0) -- (2.5,0) -- (3,1) -- (8, 1.5);


  % tap points on "a"
  \path (3,2) node[tap] {} (3,1) node[tap] {} (4.4, 1.14) node[tap] {} (4.1,1.89) node[tap] {} (5, 1.8) node[tap] {};

  \node at (3.2,2.2) {$a$};
  \node at (3.2,0.8) {$b$};

  % common right endpoint and small “)” marker
  \path (8,1.5) node[tap] {};
  \node at (8.45,1.5) {1};

\end{circuitikz}
\label{fig1}
\end{figure}

در مدار شکل \ref{fig1} جریان را در بازه ۱۰۰ تا ۵۰۰ میلی‌آمپر تغییردادیم و اختلاف پتانسیل دو سر سیم $a$ و $b$ را در جدول \ref{tab1} گزارش کردیم.

\begin{table}[h!]
\caption{}
\centering
\begin{tabular}{cc}
\hline
\lr{I (mA)} & \lr{V (v)} \\ \hline
100 & $2.62$ \\
200 & $5.31$ \\
300 & $7.98$ \\
400 & $10.63$ \\
500 & $13.27$  \\
\hline
\end{tabular}
\label{tab1}
\end{table}


\begin{figure}[h!]
\centering
\pgfplotsset{
    width=9cm,
    height=5cm,
    scale only axis,
    tick align = outside,
    yticklabel style={/pgf/number format/fixed},  
}
\begin{tikzpicture}
\begin{axis}[
title={
جریان حسب بر را سیم سر دو پتانسیل اختلاف تغییرات نمایش منحنی
},
xlabel={
\lr{[A]}
جریان
},
ylabel={
\lr{[v]}
اختلاف پتانسیل 
},
ymin=0, ymax=15,
xmin=0, xmax=0.6,
ytick distance=2.5,
xtick distance=0.1,
legend pos=north west,
ymajorgrids=true,
xmajorgrids=true,
grid style=dashed,
grid=both,
minor y tick num=1,
minor x tick num=1,
xticklabel style={
            /pgf/number format/fixed,
            /pgf/number format/precision=2,
            /pgf/number format/fixed zerofill
        },
]

\addplot[
mark=square, only marks
]
coordinates {
(0.100 , 2.62)
(0.200 , 5.31)
(0.300 , 7.98)
(0.400 , 10.63)
(0.500 ,13.27)
};

\addplot[blue] {26.62*x - 0.024};
\addlegendentry{$y = 26.62x - 0.024$}
    
\end{axis}
\end{tikzpicture}
\end{figure}

با استفاده از رگراسیون خطی بهترین منحنی را $y = 26.62x - 0.024$ به‌دست آوردیم.

می‌دانیم که رابطه‌ی بین ولتاژ و جریان طبق قانون اهم:
\[V=RI\]
است، که اگر خطی رسم کنیم و از رگرسیون خطی استفاده کنیم:
\[V = aI + b\]
در این‌جا $a = 26.62 \, \Omega$ و $b = -0.024 \, \text{\lr{V}}$

بنابراین مقاومت سیم برابر است با:
\[R = 26.62 \, \Omega\]

\begin{table}[h!]
\caption{درصدخطا}
\centering
\begin{tabular}{cccc}
\hline
\lr{$I$ (mA)} & \lr{$V$ (V)} & \lr{$R = \frac{V}{I}$ ($\Omega$)} & \lr{Error (\%)} \\ \hline
$100$ & $2.62$ & $26.20$ & $1.6$ \\
$200$ & $5.31$ & $26.55$ & $0.26$ \\
$300$ & $7.98$ & $26.60$ & $0.08$ \\
$400$ & $10.63$ & $26.58$ & $0.15$ \\
$500$ & $13.27$ & $26.54$ & $0.30$ \\
\hline
\end{tabular}
\end{table}

می‌بینی که خطاها کمتر از ۲٪ هستند؛ پس داده‌ها کاملاً با مدل خطی و مقدار $R = 26.62 \, \Omega$ سازگارند.

ضریب عرض از مبدأ 
$b = -0.024$
 ولت است که مقدار خیلی کوچکی است (در حد خطای ابزار اندازه‌گیری).
بنابراین به‌صورت عملی می‌توان گفت خط از مبدأ می‌گذرد، چون طبق قانون اهم برای یک رسانای اهمی ایده‌آل، وقتی 
$I = 0$
، باید 
$V=0$ باشد. انحراف جزئی ناشی از خطای اندازه‌گیری یا مقاومت تماس اتصالات است.

چون رابطه‌ی بین 
$V$ و 
$I$ خطی است (خط مستقیم با شیب ثابت)، و خط تقریباً از مبدأ عبور می‌کند، پس سیم از نوع اهمی (\lr{Ohmic}) است؛ یعنی از قانون اهم پیروی می‌کند.

\section{ﺑﺴﺘﮕﯽ  ﻣﻘﺎوﻣﺖ اﻟﮑﺘﺮﯾﮑﯽ ﺑﻪ ﻃﻮل ﺳﯿﻢ $R=f(L)$}

\begin{figure}[h!]
\centering
\caption{مدار ساده اندازه‌گیری}
\begin{circuitikz}
  %--- battery on the left ---
  \draw (0,0) to[battery1] (0,3);
  % polarity labels
  %\node at (-0.7,2.3) {$-$};
  %\node at (-0.7,0.7) {$+$};

  % top and bottom rails to the right
  \draw (0,3) -- (2,3);
  \draw (0,0) -- (1,0) to[ammeter] (2,0);

  % vertical voltmeter branch
  \draw (2.5,3) to[voltmeter] (5,3) -- (5, 1.8);



  %--- diverging leads a and b with tap points (red dots) ---
  % style for tap points
  \tikzset{tap/.style={draw=red!60!black, fill=red, circle, inner sep=1.2pt}}

  % upper lead "a"
  \draw
   	(2,3) -- (2.5,3) -- (3,2) -- (8,1.5);
  \draw
   	(2,0) -- (2.5,0) -- (3,1) -- (8, 1.5);


  % tap points on "a"
  \path (3,2) node[tap] {} (3,1) node[tap] {} (4.4, 1.14) node[tap] {} (4.1,1.89) node[tap] {} (5, 1.8) node[tap] {};

  \node at (3.2,2.2) {$a$};
  \node at (3.2,0.8) {$b$};
  \node at (5.7,2) {27 \lr{cm}};
  \node at (4.05,2.2) {10 \lr{cm}};
  \node at (4.4,0.8) {80 \lr{cm}};
  \node at (8,1.1) {100 \lr{cm}};

  % common right endpoint and small “)” marker
  \path (8,1.5) node[tap] {};
  \node at (8.45,1.5) {1};

\draw[->]
  (5,3) -- (5, 1.8);

\end{circuitikz}
\label{fig2}
\end{figure}

با استفاده از ولت‌متر برای طول‌های داده شده در جدول \ref{tab2}، اختلاف پتانسیل را نسبت به نقطه $a$ اندازه‌گیری کردیم و در جدول \ref{tab2} گزارش کردیم.

\begin{table}[h!]
\caption{}
\centering
\begin{tabular}{ccc}
\hline
\lr{$l$ (cm)} & \lr{V (v)} & \lr{R ($\Omega$)} \\ \hline
10  & $0.68$ & $2.72$\\
27  & $1.79$ & $7.16$\\
50  & $3.30$ & $13.21$\\
80  & $5.32$ & $21.23$\\
100 & $6.61$ & $26.44$\\
\hline
\multicolumn{3}{c}{\lr{I = $250$ (mA)}}
\end{tabular}
\label{tab2}
\end{table}

\begin{figure}[h!]
\centering
\pgfplotsset{
    width=9cm,
    height=5cm,
    scale only axis,
    tick align = outside,
    yticklabel style={/pgf/number format/fixed},  
}
\begin{tikzpicture}
\begin{axis}[
title={
$l$ طول به نسبت $R$ تغییرات نمایش منحنی
},
xlabel={
\lr{[cm]}
طول
},
ylabel={
\lr{[$\Omega$]}
مقاومت
},
ymin=0, ymax=30,
xmin=0, xmax=1,
ytick distance=5,
xtick distance=0.1,
legend pos=north west,
ymajorgrids=true,
xmajorgrids=true,
grid style=dashed,
grid=both,
minor y tick num=1,
minor x tick num=1,
xticklabel style={
            /pgf/number format/fixed,
            /pgf/number format/precision=2,
            /pgf/number format/fixed zerofill
        },
]

\addplot[
mark=square, only marks
]
coordinates {
(0.10 ,2.72)
(0.27 ,7.16)
(0.50 ,13.21)
(0.80 ,21.23)
(1.00,26.44)
};

\addplot[blue] {26.40*x+0.05};
\addlegendentry{$y=26.40x+0.05$}
    
\end{axis}
\end{tikzpicture}
\end{figure}

با استفاده از رگراسیون خطی بهترین محنی $y=26.40x+0.05$ به دست‌ آمد که یعنی شیب خط برابر  $26.40$ است.

\section{ﺗﺎﺑﻌﯿﺖ ﻣﻘﺎوﻣﺖ ﺑﺎ ﻗﻄﺮ ﺳﯿﻢ $R=f(s)$}

ابتدا سطح مقطع سیم‌ها را به دست آورد وو سپس منحنی نمایش تغییرات مقاومت برحسب عکس سطح مقطع سیم را نمایش می‌دهیم.

\begin{table}[h!]
\caption{}
\centering
\begin{tabular}{cccc}
\hline
شماره سیم & \lr{(mm)} قطر & \lr{V (v)} & \lr{R ($\Omega$)} \\
\hline
$(1) \; a,b$ & $0.25$ & $6.62$ & $26.48$ \\
$(2) \; c,d$ & $0.40$ & $2.13$ & $8.52$ \\
$(3) \; e,f$ & $0.30$ & $3.80$ & $15.20$ \\
\hline
\multicolumn{4}{c}{\lr{I = $250$ (mA)}}
\end{tabular}
\end{table}


$$
\begin{aligned}
a,b) & (1.25 \times 10^{-4})^2 \times \pi=4.9 \times 10^{-8}\\
c,d) &(2 \times 10^{-4})^2 \times \pi=1.26 \times 10^{-7}\\
e,f) &(1.5 \times 10^{-4})^2 \times \pi=7.06 \times 10^{-8}\\
\end{aligned}
$$


حال هر دو داده را داریم و از روش کمترین مربعات شیب خط را به دست می‌آوریم که مقدار آن طبق فرمول برابر

$$R=0.0348-0.4500$$

حال منحنی تتغییرات مقاومت برحسب عکس سطح مقطع سیم را داریم:

\begin{table}[h!]
\caption{}
\centering
\begin{tabular}{cc}
\hline
\lr{R ($\Omega$)}
& \lr{$\frac{1}{s}$ (mm$^{-2}$)}  \\
\hline
$26.48$ & $20.36$\\
$15.20$& $56.49$ \\
$8.52$ & $126.60$ \\
\hline
\end{tabular}
\end{table}


\begin{figure}[h!]
\centering
\pgfplotsset{
    width=9cm,
    height=5cm,
    scale only axis,
    tick align = outside,
    yticklabel style={/pgf/number format/fixed},  
}
\begin{tikzpicture}
\begin{axis}[
title={
مقاومت با سطح ارتباط
},
xlabel={
\lr{[$\Omega$]}
مقاومت
},
ylabel={
\lr{[mm$^{-2}$]}
سطح
},
ymin=0, ymax=130,
xmin=0, xmax=30,
ytick distance=20,
xtick distance=5,
legend pos=north west,
ymajorgrids=true,
xmajorgrids=true,
grid style=dashed,
grid=both,
minor y tick num=1,
minor x tick num=1,
xticklabel style={
            /pgf/number format/fixed,
            /pgf/number format/precision=2,
            /pgf/number format/fixed zerofill
        },
]

\addplot[
mark=square, only marks
]
coordinates {
(26.48, 126.60)
(15.20, 56.49)
(8.52, 20.36)
};

\addplot[blue] {0.0348*x -0.45};
%\addlegendentry{$y=26.40x+0.05$}
    
\end{axis}
\end{tikzpicture}
\end{figure}

\section{ﺗﺎﺑﻌﯿﺖ ﻣﻘﺎوﻣﺖ ﺑﺎ ﻣﻘﺎوﻣﺖ وﯾﮋه $R=f(\rho)$}

\begin{table}[h!]
\caption{}
\centering
\begin{tabular}{rcc}
\hline
جنس و شماره سیم & \lr{V (v)} & \lr{R ($\Omega$)} \\ \hline
$(3) \; e,f$ کروم نیکل & $3.81$ & $15.24$\\
$(4) \; g,h$ گالوانیزه & $0.86$ & $3.44$ \\
$(5) \; i,j$ کروم خالص & $3.00$ & $12.00$\\
\hline
\multicolumn{3}{c}{\lr{I = $250$ (mA)}}
\end{tabular}

\end{table}

مقدار مقاومت وسژه را برحسب شیب نمودار قبل مخاسبه می‌کنیم:
\[R=\rho \frac{l}{s} \Rightarrow \frac{R}{l} = \frac{\rho}{s}\]

پس با توجه به این‌که شیب نمودار $\frac{\rho}{s}$ است و مقدار مقاومت مساحت $0.0491$ \lr{mm $^2$} است.
\[\rho_1 = \frac{Rs}{l} = 0.02696 \times 10^{-5} \text{\lr{$\Omega$ m}}\]

مقدار مقاومت ویژه با توجه به نمودار(طول سیم‌ها برابر ۱ متر است.):
\[R =\rho \frac{l}{s} \Rightarrow \frac{R}{\frac{1}{s}} =\rho l\]

\[\rho_2 = \frac{Rs}{l} = 0.0346 \times 10^{-5} \text{\lr{$\Omega$ m}}\]

حال مقدار متوسط مقاومت ویژه:
\[\bar{\rho} = \frac{\rho_1 + \rho_2}{2} = 0.033 \times 10{-5} \text{\lr{$\Omega$ m}}\]

تابعیت مقاومت با مقاومت ویژه را $R = f(\rho)$ را نشان می‌دهیم و مقاومت ویژه را با توحه به رابطه $R = \rho \frac{l}{A}$ به دست می‌اوریم:

\[
\begin{aligned}
R_3 & = \frac{V_3}{I} = 3.91 \Omega \rightarrow \rho_3 = 12.27 \times 10 ^{-8} \Omega \text{\lr{m}} \\
R_3 & = \frac{V_3}{I} = 35.48 \Omega \rightarrow \rho_3 = 111.45 \times 10 ^{-8} \Omega \text{\lr{m}} \\
R_3 & = \frac{V_3}{I} = 18.90 \Omega \rightarrow \rho_3 = 59.47 \times 10 ^{-8} \Omega \text{\lr{m}} \\
\end{aligned}
\]
\end{document}